\subsection{Dual of a reflexive Fréchet space}

PROPOSITION 4. --- Let E be a reflexive Fréchet space. The strong dual \(E_{b}^{\prime}\) of E is the inductive limit of a sequence of Banach spaces.

Let \((\mathrm{V}_{n})_{n\in\mathbf{N}}\) be a decreasing sequence of convex, balanced and closed neighbourhoods of 0 in E, such that every neighbourhood of 0 in E contains one of the \(V_{n}\) . Let \(A_{n}\) be the polar of \(V_{n}\) in \(E'\) . Then \(A_{n}\) is convex, balanced and compact for \(\sigma(\mathrm{E}',\mathrm{E})\) ; by III, p.~8, corollary the space \(E_{A_{n}}'\) is a Banach space. We shall prove that \(E_{b}'\) is the inductive limit of the spaces \(E_{A_{n}}'\) ; in other words, that every convex and balanced subset U of \(E'\) which absorbs each of the \(A_{n}\) is a neighbourhood of 0 in \(E_{b}'\) . For every \(n\in\mathbf{N}\) , choose a real number \(\lambda_{n}>0\) such that \(\lambda_{n}A_{n}\subset U\) . Let \(B_{n}\) be the convex envelope of the set \(\bigcup_{0\leqslant i\leqslant n}\lambda_{i}A_{i}\) ; put \(V=\bigcup_{n}B_{n}\) , then \(V\subset U\) . For every \(n\in\mathbf{N}\) , the set \(B_{n}\) is convex, balanced and compact for \(\sigma(\mathrm{E}',\mathrm{E})\) (II, p.~14, prop. 15).

Now we shall show that \(\frac{1}{2}\mathrm{V}^{\circ \circ}\subset \mathrm{V}\) . Let \(x\in \mathrm{E}_b^\prime -\mathrm{V}\) ; for every \(n\in \mathbf{N}\) , we have \(x\notin \mathbf{B}_n\) , and since \(\mathbf{B}_n\) is closed for \(\sigma (\mathrm{E}',\mathrm{E})\) there exists an element \(y_{n}\) in \(\mathbf{B}_n^\circ\) such that \(\langle y_n, x \rangle = 1\) (II, p.~38, prop. 4). Since E is reflexive, every bounded subset of E is relatively compact for \(\sigma(E, E')\) (IV, p.~16, th. 2). By the definition of \(B_n\) , we have

\[
\lambda_ {i} y _ {n} \in \mathrm{V} _ {i} \quad \text { for   all } \quad n \geqslant i  ,\tag{3}
\]

hence the sequence \((y_{n})\) is bounded. Let \(y\) be a limit point of \((y_{n})\) for the topology \(\sigma(\mathrm{E},\mathrm{E}')\) . We have \(y\in\mathrm{V}^{\circ}=\bigcap_{n}\mathrm{B}_{n}^{\circ}\) and \(\langle y,x\rangle=1\) . Hence \(x\notin\frac{1}{2}\mathrm{V}^{\circ\circ}\) , and so we have the inclusion \(\frac{1}{2}\mathrm{V}^{\circ\circ}\subset\mathrm{V}\) and a fortiori, \(\frac{1}{2}\mathrm{V}^{\circ\circ}\subset\mathrm{U}\) .

Since every bounded subset of \(E_{b}^{\prime}\) is contained in one of the sets \(A_{n}\) , the set \(V = \bigcup_{n} B_{n}\)

absorbs every bounded subset of \(E_{b}^{\prime}\) . Consequently, \(V^{\circ}\) is bounded in E, hence \(\frac{1}{2}V^{\circ\circ}\) is a neighbourhood of 0 in \(E_{b}^{\prime}\) . A fortiori, U is a neighbourhood of 0 in \(E_{b}^{\prime}\) .

COROLLARY. --- The strong dual of a reflexive Fréchet space is bornological and barrelled.

An inductive limit of Banach spaces is bornological by definition. Further, a Banach space is barrelled (III, p.~25, corollary) and every inductive limit of barrelled spaces is a barrelled space (III, p.~25, cor. 3).

